
% fala um pouco sobre as condições iniciais da tabela de páginas:
% paginação não ativada, mas virtual mapeado diretamente para físico,
% que é como as coisas ficam quando ativamos a paginação também
% já que a paginação é ativada após criarmos o primeiro processo.
% menciona por que ainda temos SEG_UCODE/SEG_UDATA?
% alguma vez dissemos realmente o que os dois bits inferiores de %cs fazem?
% em particular, sua interação com PTE_U
% barra lateral sobre por que é extern char[]
\chapter{Tabelas de páginas}
\label{CH:MEM}

Tabelas de páginas são o mecanismo mais popular por meio do qual o sistema operacional fornece a cada processo seu próprio espaço de endereçamento e memória privados. Tabelas de páginas determinam 
o que significam os endereços de memória e quais partes da memória física podem ser acessadas. Elas permitem que o xv6 isole os espaços de endereçamento de diferentes processos e os multiplexe em 
uma única memória física. Tabelas de páginas são um projeto popular porque fornecem um nível de indireção que permite aos sistemas operacionais realizar muitos truques. O xv6 realiza alguns 
truques: mapeando a mesma memória (uma página de trampolim) em vários espaços de endereçamento e protegendo pilhas de kernel e de usuário com uma página não mapeada. O restante deste capítulo 
explica as tabelas de páginas fornecidas pelo hardware RISC-V e como o xv6 as utiliza.

%%
\section{Hardware de paginação}
%%

Como lembrança, instruções RISC-V (tanto de usuário quanto de kernel) manipulam endereços virtuais. A RAM da máquina, ou memória física, é indexada com endereços físicos. O hardware de tabela de 
páginas do RISC-V conecta esses dois tipos de endereços, mapeando cada endereço virtual para um endereço físico.

\begin{figure}[t]
\center
\includegraphics[scale=0.5]{fig/riscv_address.pdf}
\caption{Endereços virtuais e físicos no RISC-V, com uma tabela de páginas lógica simplificada.}
\label{fig:riscv_address}
\end{figure}

O xv6 roda no RISC-V Sv39, o que significa que apenas os 39 bits inferiores de um endereço virtual de 64 bits são usados; os 25 bits superiores não são utilizados. Nessa configuração Sv39, uma 
tabela de páginas RISC-V é logicamente um array de $2^{27}$ (134.217.728) \indextext{entradas de tabela de páginas (PTEs)}. Cada PTE contém um número de página física (PPN) de 44 bits e alguns 
flags. O hardware de paginação traduz um endereço virtual usando os 27 bits superiores dos 39 bits para indexar na tabela de páginas e encontrar uma PTE, e gerando um endereço físico de 56 bits 
cujos 44 bits superiores vêm do PPN da PTE e cujos 12 bits inferiores são copiados do endereço virtual original. A Figura~\ref{fig:riscv_address} mostra esse processo com uma visão lógica da tabela 
de páginas como um array simples de PTEs (veja a Figura~\ref{fig:riscv_pagetable} para uma explicação mais completa). Uma tabela de páginas dá ao sistema operacional controle sobre traduções de 
endereços virtuais para físicos na granularidade de blocos alinhados de 4096 ($2^{12}$) bytes. Tal bloco é chamado de \indextext{página}.

\begin{figure}[t]
\center
\includegraphics[scale=0.5]{fig/riscv_pagetable.pdf}
\caption{Detalhes da tradução de endereços no RISC-V.}
\label{fig:riscv_pagetable}
\end{figure}

No RISC-V Sv39, os 25 bits superiores de um endereço virtual não são usados para tradução. O endereço físico também tem espaço para expansão: há espaço no formato da PTE para que o número da página 
física cresça mais 10 bits. Os projetistas do RISC-V escolheram esses números com base em previsões tecnológicas. $2^{39}$ bytes é 512 GB, o que deve ser espaço suficiente para aplicações rodando 
em computadores RISC-V. $2^{56}$ é espaço físico suficiente para o futuro próximo, acomodando muitos dispositivos de E/S e chips de RAM. Se for necessário mais, os projetistas do RISC-V definiram o 
Sv48 com endereços virtuais de 48 bits~\cite{riscv:priv}.

Como mostra a Figura~\ref{fig:riscv_pagetable}, uma CPU R_

%%
\section{Espaço de endereçamento do kernel}
%%

O xv6 mantém uma tabela de páginas por processo, descrevendo o espaço de endereçamento de usuário de cada processo, além de uma única tabela de páginas que descreve o espaço de endereçamento do 
kernel. O kernel configura o layout de seu espaço de endereçamento para ter acesso à memória física e a vários recursos de hardware em endereços virtuais previsíveis. A Figura~\ref{fig:xv6_layout} 
mostra como esse layout mapeia os endereços virtuais do kernel para endereços físicos. O arquivo \fileref{kernel/memlayout.h} declara as constantes para o layout de memória do kernel no xv6.

O QEMU simula um computador que inclui RAM (memória física) começando no endereço físico \texttt{0x80000000} e se estendendo até pelo menos \texttt{0x88000000}, que o xv6 chama de \texttt{PHYSTOP}. 
A simulação do QEMU também inclui dispositivos de E/S, como uma interface de disco. O QEMU expõe as interfaces dos dispositivos para o software como registradores de controle \indextext{mapeados em 
memória}, que ficam abaixo de \texttt{0x80000000} no espaço de endereçamento físico. O kernel pode interagir com os dispositivos lendo/escrevendo nesses endereços físicos especiais; tais leituras e 
escritas se comunicam com o hardware do dispositivo, e não com a RAM. O Capítulo~\ref{CH:TRAP} explica como o xv6 interage com os dispositivos.

O kernel acessa a RAM e os registradores mapeados em memória dos dispositivos usando o “mapeamento direto”; isto é, mapeando os recursos em endereços virtuais iguais ao endereço físico. Por 
exemplo, o próprio kernel está localizado em \lstinline{KERNBASE=0x80000000} tanto no espaço de endereçamento virtual quanto na memória física. O mapeamento direto simplifica o código do kernel que 
lê ou escreve memória física. Por exemplo, quando \lstinline{fork} aloca memória de usuário para o processo filho, o alocador retorna o endereço físico dessa memória; \lstinline{fork} usa esse 
endereço diretamente como um endereço virtual ao copiar a memória de usuário do processo pai para o filho.

Há alguns endereços virtuais do kernel que não são mapeados diretamente:

\begin{itemize}

\item A página trampolim. Ela é mapeada no topo do espaço de endereçamento virtual; as tabelas de páginas de usuário também possuem esse mesmo mapeamento. O Capítulo~\ref{CH:TRAP} discute o papel 
da página trampolim, mas aqui vemos um caso interessante de uso de tabelas de páginas: uma página física (contendo o código do trampolim) é mapeada duas vezes no espaço de endereçamento virtual do 
kernel — uma vez no topo do espaço de endereçamento virtual e outra vez com mapeamento direto.

\item As páginas de pilha do kernel. Cada processo tem sua própria pilha de kernel, que é mapeada em endereços altos, de modo que abaixo dela o xv6 pode deixar uma \indextext{página de proteção} 
não mapeada. A PTE da página de proteção é inválida (ou seja, \lstinline{PTE_V} não está definida), de modo que, se o kernel ultrapassar os limites da pilha de kernel, provavelmente causará uma 
exceção e o kernel entrará em pânico. Sem uma página de proteção, uma pilha excedente sobrescreveria outras partes da memória do kernel, resultando em funcionamento incorreto. Um travamento com 
pânico é preferível.

\end{itemize}

Embora o kernel use suas pilhas por meio dos mapeamentos em memória alta, elas também estão acessíveis ao kernel por meio de um endereço com mapeamento direto. Um projeto alternativo poderia usar 
apenas o mapeamento direto, acessando as pilhas nesse endereço diretamente mapeado. No entanto, nessa configuração, fornecer páginas de proteção exigiria o desmapeamento de endereços virtuais que, 
de outra forma, se refeririam à memória física, o que dificultaria seu uso.

O kernel mapeia as páginas do trampolim e do código do kernel com as permissões \lstinline{PTE_R} e \lstinline{PTE_X}. O kernel lê e executa instruções a partir dessas páginas. O kernel mapeia as 
outras páginas com as permissões \lstinline{PTE_R} e \lstinline{PTE_W}, para que possa ler e escrever a memória nessas páginas. Os mapeamentos para as páginas de proteção são inválidos.

%% 
\section{Código: criação de um espaço de endereçamento}
%% 

A maior parte do código do xv6 para manipular espaços de endereçamento e
tabelas de páginas reside em {\tt vm.c}
\lineref{kernel/vm.c:1}.
A estrutura de dados central é {\tt pagetable\_t}, que
é na verdade um ponteiro para uma página de tabela de páginas raiz do RISC-V;
um {\tt pagetable\_t} pode ser tanto a tabela de páginas do kernel
quanto uma das tabelas de páginas por processo.
As funções centrais são {\tt walk},
que encontra a PTE para um endereço virtual,
e {\tt mappages}, que instala PTEs para novos mapeamentos.
Funções que começam com {\tt kvm} manipulam
a tabela de páginas do kernel; funções que começam com {\tt uvm}
manipulam uma tabela de páginas de usuário; outras funções são
usadas para ambas.
{\tt copyout} e {\tt copyin} copiam dados de e para
endereços virtuais de usuário fornecidos como argumentos de chamadas de sistema;
elas estão em {\tt vm.c} porque precisam traduzir explicitamente esses endereços
para encontrar a memória física correspondente.

Logo no início da sequência de boot,
\indexcode{main}
chama
\indexcode{kvminit}
\lineref{kernel/vm.c:/^kvminit/}
para criar a tabela de páginas do kernel usando
\indexcode{kvmmake}
\lineref{kernel/vm.c:/^kvmmake/}.
Essa chamada ocorre antes que o xv6 habilite paginação no RISC-V,
portanto os endereços referem-se diretamente à memória física.
\lstinline{kvmmake}
primeiro aloca uma página de memória física para conter a página de tabela de páginas raiz.
Depois, chama
\indexcode{kvmmap}
para instalar as traduções que o kernel precisa.
As traduções incluem as instruções e dados do kernel, memória física até
\indexcode{PHYSTOP},
e intervalos de memória que na verdade são dispositivos.
\indexcode{proc_mapstacks}
\lineref{kernel/proc.c:/^proc_mapstacks/}
aloca uma pilha de kernel para cada
processo. Ela chama \lstinline{kvmmap} para mapear cada pilha no endereço virtual gerado por
\lstinline{KSTACK}, o que deixa espaço para as páginas de proteção de pilha inválidas.

\indexcode{kvmmap} 
\lineref{kernel/vm.c:/^kvmmap/}
chama
\indexcode{mappages}
\lineref{kernel/vm.c:/^mappages/},
que
instala mapeamentos em uma tabela de páginas
para um intervalo de endereços virtuais para
um intervalo correspondente de endereços físicos.
Faz isso separadamente para cada endereço virtual no intervalo,
em passos de página.
Para cada endereço virtual a ser mapeado,
\lstinline{mappages}
chama
\indexcode{walk}
para encontrar o endereço da PTE correspondente.
Em seguida, inicializa a PTE para conter o número da página física relevante,
as permissões desejadas 
(\lstinline{PTE_W},
\lstinline{PTE_X},
e/ou
\lstinline{PTE_R}),
e
\lstinline{PTE_V}
para marcar a PTE como válida
\lineref{kernel/vm.c:/perm...PTE_V/}.

\indexcode{walk}
\lineref{kernel/vm.c:/^walk/}
imita o hardware de paginação do RISC-V ao
procurar a PTE de um endereço virtual (ver
Figura~\ref{fig:riscv_pagetable}).
\lstinline{walk}
desce a tabela de páginas um nível por vez,
usando os 9 bits de endereço virtual de cada nível para
indexar a página de diretório correspondente.
Em cada nível, encontra-se a PTE da
página de diretório do próximo nível, ou a PTE da
página final
\lineref{kernel/vm.c:/pte.=..pagetable/}.
Se uma PTE em um diretório de primeiro ou segundo nível
não for válida, então a página de diretório necessária
ainda não foi alocada;
se o argumento
\lstinline{alloc}
estiver ativado,
\lstinline{walk}
aloca uma nova página de tabela de páginas e armazena seu endereço físico na PTE.
Ela retorna o endereço da PTE no nível mais baixo da árvore
\lineref{kernel/vm.c:/return..pagetable/}.

O código acima depende de a memória física estar mapeada diretamente no
espaço de endereçamento virtual do kernel. Por exemplo, à medida que \lstinline{walk} desce os níveis
da tabela de páginas, ela extrai o endereço (físico) da
próxima tabela de páginas a partir de uma PTE \lineref{kernel/vm.c:/pagetable.=..pa.*E2P/},
e então usa esse endereço como um
endereço virtual para buscar a PTE no próximo nível
\lineref{kernel/vm.c:/t..pte.=..paget/}.

\indexcode{main}
chama
\indexcode{kvminithart}
\lineref{kernel/vm.c:/^kvminithart/}
para instalar a tabela de páginas do kernel.
Ela escreve o endereço físico da página de tabela de páginas raiz
no registrador
\texttt{satp}.
Depois disso, a CPU passará a traduzir endereços usando a tabela
de páginas do kernel. Como o kernel usa um mapeamento direto, o
endereço virtual da próxima instrução agora será mapeado corretamente para a
memória física correspondente.

Cada CPU RISC-V armazena em cache entradas da tabela de páginas em um
\indextext{Translation Look-aside Buffer (TLB)}, e quando o xv6 altera
uma tabela de páginas, ele deve informar à CPU para invalidar as entradas
correspondentes em cache no TLB. Caso contrário,
em algum momento posterior o TLB pode
usar um mapeamento antigo em cache, apontando para uma página física que, nesse meio tempo,
tenha sido alocada para outro processo, e como resultado, um processo
poderia sobrescrever a memória de outro processo. O RISC-V
possui uma instrução \indexcode{sfence.vma} que limpa
o TLB da CPU atual.
O xv6 executa {\tt sfence.vma} em {\tt kvminithart} após recarregar o 
registrador \texttt{satp}, e também no código trampoline que
troca para a tabela de páginas de usuário antes de retornar ao espaço do usuário
\lineref{kernel/trampoline.S:/sfence.vma/}.

Também é necessário emitir \texttt{sfence.vma} antes de alterar
\texttt{satp}, a fim de esperar pela conclusão de todas as operações de leitura e
escrita pendentes. Essa espera garante que as atualizações anteriores à tabela de páginas tenham
sido concluídas e que as leituras e escritas anteriores usem a tabela de páginas antiga,
não a nova.

Para evitar limpar completamente o TLB, CPUs RISC-V podem suportar
identificadores de espaço de endereçamento (ASIDs)~\cite{riscv:priv}. O kernel pode então
limpar apenas as entradas do TLB para um espaço de endereçamento específico. O xv6
não utiliza esse recurso.

\section{Alocação de memória física}

O kernel deve alocar e liberar memória física em tempo de execução para
tabelas de páginas,
memória de usuário,
pilhas do kernel
e buffers de pipe.

O xv6 usa a memória física entre o final do kernel e
\indexcode{PHYSTOP}
para alocação em tempo de execução. Ele aloca e libera páginas inteiras de 4096 bytes
por vez. Ele rastreia quais páginas estão livres encadeando uma
lista ligada através das próprias páginas. A alocação consiste em
remover uma página da lista ligada; liberar consiste em adicionar a
página liberada de volta à lista.

%%
\section{Código: Alocador de memória física}
%%

O alocador está localizado em {\tt kalloc.c} \lineref{kernel/kalloc.c:1}.
A estrutura de dados do alocador é uma
\textit{lista livre}
de páginas de memória física disponíveis
para alocação.
Cada elemento da lista correspondente a uma página livre é uma
\indexcode{struct run}
\lineref{kernel/kalloc.c:/^struct.run/}.
De onde o alocador obtém a memória
para armazenar essa estrutura de dados?
Ele armazena a estrutura
\lstinline{run}
de cada página livre na própria página,
já que não há mais nada armazenado ali.
A lista livre é
protegida por uma trava do tipo spin-lock
\linerefs{kernel/kalloc.c:/^struct.{/,/}/}.
A lista e a trava são empacotadas em uma \textit{struct}
para deixar claro que a trava protege os campos
dessa estrutura.
Por ora, ignore a trava e as chamadas a
\lstinline{acquire}
e
\lstinline{release};
o Capítulo~\ref{CH:LOCK} examinará
travas em detalhe.

A função
\indexcode{main}
chama
\indexcode{kinit}
para inicializar o alocador
\lineref{kernel/kalloc.c:/^kinit/}.
\lstinline{kinit}
inicializa a lista livre para conter
todas as páginas entre o fim do kernel e {\tt PHYSTOP}.
O xv6 deveria determinar quanta memória física
está disponível analisando as informações de configuração
fornecidas pelo hardware.
Em vez disso, o xv6 assume que a máquina possui
128 megabytes de RAM.
\lstinline{kinit}
chama
\indexcode{freerange}
para adicionar memória à lista livre por meio de chamadas por página à
\indexcode{kfree}.
Uma PTE só pode referenciar um endereço físico que esteja alinhado
em um limite de 4096 bytes (ou seja, seja múltiplo de 4096), então
\lstinline{freerange}
usa
\indexcode{PGROUNDUP}
para garantir que libera apenas endereços físicos alinhados.
O alocador começa sem nenhuma memória;
essas chamadas à
\lstinline{kfree}
lhe fornecem memória para gerenciar.

O alocador às vezes trata endereços como inteiros
para realizar operações aritméticas sobre eles (por exemplo,
percorrer todas as páginas em
\lstinline{freerange}),
e às vezes usa endereços como ponteiros para ler e
escrever memória (por exemplo, manipulando a
estrutura
\lstinline{run}
armazenada em cada página);
esse uso duplo de endereços é a principal razão pela qual o
código do alocador está repleto de conversões de tipo em C.
\index{conversão de tipo}

A função
\lstinline{kfree}
\lineref{kernel/kalloc.c:/^kfree/}
começa definindo cada byte da
memória sendo liberada com o valor 1.
Isso fará com que código que use memória após ela ter sido liberada
(utilize “referências pendentes”)
leia lixo em vez dos antigos conteúdos válidos;
com sorte, isso fará com que tal código falhe mais rapidamente.
Em seguida,
\lstinline{kfree}
adiciona a página ao início da lista livre:
ela converte
\lstinline{pa}
para um ponteiro para
\lstinline{struct}
\lstinline{run},
registra o antigo início da lista livre em
\lstinline{r->next},
e define a lista livre como sendo
\lstinline{r}.
\indexcode{kalloc}
remove e retorna o primeiro elemento da lista livre.

\section{Espaço de endereçamento do processo}

Cada processo possui sua própria tabela de páginas e, quando o xv6 troca de processo, ele também altera a tabela de páginas.  
A Figura~\ref{fig:processlayout} mostra o espaço de endereçamento de um processo com mais detalhes do que a Figura~\ref{fig:as}.  
A memória do usuário de um processo começa no endereço virtual zero e pode crescer até \texttt{MAXVA}
\lineref{kernel/riscv.h:/MAXVA/}, permitindo que um processo enderece, em princípio, 256 Gigabytes de memória.

O espaço de endereçamento de um processo consiste em páginas que contêm o texto do programa (que o xv6 mapeia com as permissões \lstinline{PTE_R}, \lstinline{PTE_X} e \lstinline{PTE_U}), páginas que contêm os dados pré-inicializados do programa, uma página para a pilha (stack) e páginas para o heap.  
O xv6 mapeia os dados, a pilha e o heap com as permissões \lstinline{PTE_R}, \lstinline{PTE_W} e \lstinline{PTE_U}.

O uso de permissões dentro de um espaço de endereçamento de usuário é uma técnica comum para fortalecer a segurança de um processo.  
Se o texto fosse mapeado com \lstinline{PTE_W}, então um processo poderia acidentalmente modificar seu próprio programa; por exemplo, um erro de programação pode fazer com que o programa escreva em 
um ponteiro nulo, modificando instruções no endereço 0, e depois continue executando, talvez causando ainda mais problemas.  
Para detectar tais erros imediatamente, o xv6 mapeia o texto sem \lstinline{PTE_W}; se um programa tentar armazenar acidentalmente no endereço 0, o hardware se recusará a executar a operação de 
escrita e gerará uma falha de página (veja Seção~\ref{sec:pagefaults}).  
O kernel então finaliza o processo e imprime uma mensagem informativa para que o desenvolvedor possa rastrear o problema.

De forma similar, ao mapear dados sem \lstinline{PTE_X}, um programa de usuário não pode acidentalmente saltar para um endereço nos dados do programa e começar a executá-los.

No mundo real, reforçar um processo ao definir permissões cuidadosamente também ajuda a defender contra ataques de segurança.  
Um adversário pode fornecer entradas cuidadosamente construídas para um programa (por exemplo, um servidor Web) que desencadeiem um bug no programa, na esperança de transformar esse bug em um 
exploit~\cite{aleph:smashing}.  
Definir permissões com cuidado e outras técnicas, como a randomização do layout do espaço de endereçamento do usuário, tornam tais ataques mais difíceis.

A pilha é uma única página, e é mostrada com o conteúdo inicial criado pela chamada \texttt{exec}.  
Strings contendo os argumentos da linha de comando, bem como um array de ponteiros para esses argumentos, estão no topo da pilha.  
Logo abaixo estão os valores que permitem que um programa comece na função \lstinline{main}, como se \lstinline{main(argc}, \lstinline{argv)} tivesse acabado de ser chamada.

Para detectar um estouro da pilha do usuário além da memória alocada, o xv6 coloca uma página de proteção inacessível logo abaixo da pilha, limpando o bit \lstinline{PTE_U}.  
Se a pilha do usuário estourar e o processo tentar usar um endereço abaixo da pilha, o hardware gerará uma exceção de falha de página, pois a página de proteção é inacessível a um programa 
executando em modo de usuário.  
Um sistema operacional real pode, em vez disso, alocar automaticamente mais memória para a pilha do usuário quando ela estourar.

Quando um processo solicita mais memória de usuário ao xv6, o xv6 aumenta o heap do processo.  
Primeiro, o xv6 usa {\tt kalloc} para alocar páginas físicas.  
Depois, ele adiciona PTEs à tabela de páginas do processo que apontam para as novas páginas físicas.  
O xv6 define os bits \lstinline{PTE_W}, \lstinline{PTE_R}, \lstinline{PTE_U} e \lstinline{PTE_V} nesses PTEs.  
A maioria dos processos não usa todo o espaço de endereçamento do usuário;  
o xv6 deixa o bit \lstinline{PTE_V} limpo nos PTEs não utilizados.

Vemos aqui alguns bons exemplos de uso das tabelas de páginas.  
Primeiro, as tabelas de páginas de diferentes processos traduzem endereços de usuário para diferentes páginas de memória física, de forma que cada processo tenha memória de usuário privada.  
Segundo, cada processo vê sua memória como tendo endereços virtuais contíguos começando em zero, enquanto a memória física do processo pode ser não contígua.  
Terceiro, o kernel mapeia uma página com código trampolim no topo do espaço de endereçamento do usuário (sem \lstinline{PTE_U}), assim uma única página de memória física aparece em todos os espaços 
de endereçamento, mas só pode ser usada pelo kernel.

\begin{figure}[t]
\center
\includegraphics[scale=0.5]{fig/processlayout.png}
\caption{Espaço de endereçamento do usuário de um processo, com sua pilha inicial.}
\label{fig:processlayout}
\end{figure}

\section{Código: sbrk}

\lstinline{sbrk} é a chamada de sistema para um processo reduzir ou expandir sua memória.  
A chamada de sistema é implementada pela função \lstinline{growproc}  
\lineref{kernel/proc.c:/^growproc/}.  
\lstinline{growproc} chama \lstinline{uvmalloc} ou \lstinline{uvmdealloc}, dependendo se \lstinline{n} é positivo ou negativo.  
\lstinline{uvmalloc}  
\lineref{kernel/vm.c:/^uvmalloc/}  
aloca memória física com {\tt kalloc}, zera a memória alocada e adiciona PTEs à tabela de páginas do usuário com {\tt mappages}.  
\lstinline{uvmdealloc} chama {\tt uvmunmap}  
\lineref{kernel/vm.c:/^uvmunmap/},  
que usa {\tt walk} para encontrar os PTEs e {\tt kfree} para liberar a memória física à qual eles se referem.

O xv6 usa a tabela de páginas de um processo não apenas para informar ao hardware como mapear os endereços virtuais de usuário, mas também como o único registro de quais páginas de memória física 
estão alocadas para aquele processo.  
Essa é a razão pela qual liberar memória de usuário (em {\tt uvmunmap}) exige a inspeção da tabela de páginas do usuário.
%% 
\section{Código: exec}
%% 
\lstinline{exec} é uma chamada de sistema que substitui o espaço de endereçamento do usuário de um processo por dados lidos de um arquivo, chamado de binário ou arquivo executável.  
Um binário é normalmente o resultado da compilação e ligação, e contém instruções de máquina e dados do programa.  
\lstinline{exec} \lineref{kernel/exec.c:/^exec/} abre o binário nomeado \lstinline{path} usando \indexcode{namei} \lineref{kernel/exec.c:/namei/}, que é explicado no Capítulo~\ref{CH:FS}.  
Em seguida, lê o cabeçalho ELF. Os binários do xv6 estão formatados no amplamente utilizado \indextext{formato ELF}, definido em \fileref{kernel/elf.h}.  
Um binário ELF consiste em um cabeçalho ELF, \indexcode{struct elfhdr} \lineref{kernel/elf.h:/^struct.elfhdr/}, seguido por uma sequência de cabeçalhos de seção de programa,  
\lstinline{struct proghdr} \lineref{kernel/elf.h:/^struct.proghdr/}.  
Cada \lstinline{proghdr} descreve uma seção da aplicação que deve ser carregada na memória; programas do xv6 têm dois cabeçalhos de seção de programa: um para instruções e um para dados.

O primeiro passo é uma verificação rápida de que o arquivo provavelmente contém um binário ELF.  
Um binário ELF começa com o número mágico de quatro bytes  
\lstinline{0x7F},  
\lstinline{`E'},  
\lstinline{`L'},  
\lstinline{`F'},  
ou  
\indexcode{ELF_MAGIC}  
\lineref{kernel/elf.h:/ELF_MAGIC/}.  
Se o cabeçalho ELF tem o número mágico correto,  
\lstinline{exec}  
assume que o binário está bem formado.

\lstinline{exec}  
aloca uma nova tabela de páginas sem mapeamentos de usuário com  
\indexcode{proc_pagetable}  
\lineref{kernel/exec.c:/proc_pagetable/},  
aloca memória para cada segmento ELF com  
\indexcode{uvmalloc}  
\lineref{kernel/exec.c:/uvmalloc/},  
e carrega cada segmento na memória com  
\indexcode{loadseg}  
\lineref{kernel/exec.c:/loadseg/}.  
\lstinline{loadseg}  
usa  
\indexcode{walkaddr}  
para encontrar o endereço físico da memória alocada onde escrever cada página do segmento ELF, e  
\indexcode{readi}  
para ler do arquivo.

O cabeçalho de seção de programa para  
\indexcode{/init},  
o primeiro programa de usuário criado com  
\lstinline{exec},  
se parece com isto:
\begin{footnotesize}
\begin{verbatim}
# objdump -p user/_init

user/_init:     file format elf64-little

Program Header:
0x70000003 off    0x0000000000006bb0 vaddr 0x0000000000000000
                                       paddr 0x0000000000000000 align 2**0
         filesz 0x000000000000004a memsz 0x0000000000000000 flags r--
    LOAD off    0x0000000000001000 vaddr 0x0000000000000000
                                       paddr 0x0000000000000000 align 2**12
         filesz 0x0000000000001000 memsz 0x0000000000001000 flags r-x
    LOAD off    0x0000000000002000 vaddr 0x0000000000001000
                                       paddr 0x0000000000001000 align 2**12
         filesz 0x0000000000000010 memsz 0x0000000000000030 flags rw-
   STACK off    0x0000000000000000 vaddr 0x0000000000000000
                                       paddr 0x0000000000000000 align 2**4
         filesz 0x0000000000000000 memsz 0x0000000000000000 flags rw-
\end{verbatim}
\end{footnotesize}

Vemos que o segmento de texto deve ser carregado no endereço virtual 0x0 na memória (sem permissões de escrita) a partir do conteúdo no deslocamento 0x1000 no arquivo.  
Também vemos que os dados devem ser carregados no endereço 0x1000, que está no limite de uma página, e sem permissões de execução.

O \lstinline{filesz} de um cabeçalho de seção de programa pode ser menor que o \lstinline{memsz},  
indicando que a diferença entre eles deve ser preenchida com zeros (para variáveis globais em C) em vez de ser lida do arquivo.  
Para \lstinline{/init}, o \lstinline{filesz} dos dados é 0x10 bytes e o \lstinline{memsz} é 0x30 bytes,  
e, portanto, \indexcode{uvmalloc} aloca memória física suficiente para armazenar 0x30 bytes, mas lê apenas 0x10 bytes do arquivo \lstinline{/init}.

Agora \indexcode{exec} aloca e inicializa a pilha do usuário.  
Ele aloca apenas uma página de pilha.  
\lstinline{exec} copia as strings de argumento para o topo da pilha uma por uma, registrando os ponteiros para elas em \indexcode{ustack}.  
Coloca um ponteiro nulo no fim da lista \indexcode{argv} que será passada para \lstinline{main}.  
Os valores de \indexcode{argc} e \lstinline{argv} são passados para \lstinline{main} através do caminho de retorno da chamada de sistema: \lstinline{argc} é passado via o valor de retorno da 
chamada de sistema, que vai em {\tt a0}, e \lstinline{argv} é passado através da entrada {\tt a1} do \textit{trapframe} do processo.

\lstinline{exec} posiciona uma página inacessível logo abaixo da página da pilha,  
de forma que programas que tentarem usar mais de uma página causarão falha.  
Essa página inacessível também permite que \lstinline{exec} lide com argumentos grandes demais;  
nessa situação, a função \indexcode{copyout} \lineref{kernel/vm.c:/^copyout/},  
usada por \lstinline{exec} para copiar argumentos para a pilha, perceberá que a página de destino não é acessível e retornará -1.

Durante a preparação da nova imagem de memória, se \lstinline{exec} detectar um erro como um segmento de programa inválido, ele salta para o rótulo \lstinline{bad},  
libera a nova imagem e retorna -1.  
\lstinline{exec} deve esperar para liberar a imagem antiga até ter certeza de que a chamada de sistema será bem-sucedida:  
se a imagem antiga for liberada, a chamada de sistema não poderá retornar -1 para ela.  
Os únicos casos de erro em \lstinline{exec} ocorrem durante a criação da imagem.  
Uma vez que a imagem esteja completa, \lstinline{exec} pode se comprometer com a nova tabela de páginas \lineref{kernel/exec.c:/pagetable.=.pagetable/}  
e liberar a antiga \lineref{kernel/exec.c:/proc_freepagetable/}.

\lstinline{exec} carrega bytes do arquivo ELF na memória nos endereços especificados pelo arquivo ELF.  
Usuários ou processos podem colocar quaisquer endereços que desejarem em um arquivo ELF.  
Assim, \lstinline{exec} é arriscado, pois os endereços no arquivo ELF podem referenciar o kernel, acidentalmente ou de propósito.  
As consequências para um kernel desavisado podem variar de uma falha até a subversão maliciosa dos mecanismos de isolamento do kernel (ou seja, uma exploração de segurança).  
O xv6 realiza várias verificações para evitar esses riscos.  
Por exemplo, \lstinline{if(ph.vaddr + ph.memsz < ph.vaddr)} verifica se a soma transborda um inteiro de 64 bits.  
O perigo é que um usuário poderia construir um binário ELF com um \lstinline{ph.vaddr} que aponta para um endereço escolhido pelo usuário,  
e um \lstinline{ph.memsz} grande o suficiente para que a soma transborde para 0x1000, que parecerá um valor válido.  
Em uma versão antiga do xv6 em que o espaço de endereçamento do usuário também continha o kernel (mas não era legível/gravável em modo usuário),  
o usuário poderia escolher um endereço que correspondesse à memória do kernel e, assim, copiar dados do binário ELF para o kernel.  
Na versão RISC-V do xv6 isso não pode acontecer, pois o kernel possui sua própria tabela de páginas separada;  
\lstinline{loadseg} carrega na tabela de páginas do processo, não na tabela do kernel.

É fácil para um desenvolvedor do kernel omitir uma verificação crucial,  
e kernels reais têm um longo histórico de verificações ausentes cuja ausência pode ser explorada por programas de usuário para obter privilégios de kernel.  
É provável que o xv6 não faça um trabalho completo de validação dos dados em nível de usuário fornecidos ao kernel,  
o que um programa de usuário malicioso poderia explorar para contornar o isolamento do xv6.

%% 
\section{Mundo real}
%% 

Como a maioria dos sistemas operacionais, o xv6 utiliza o hardware de paginação para proteção e mapeamento de memória. A maioria dos sistemas operacionais faz um uso muito mais sofisticado da 
paginação do que o xv6, combinando paginação e exceções de falha de página, o que discutiremos no Capítulo~\ref{CH:TRAP}.

O xv6 é simplificado pelo uso de um mapeamento direto entre endereços virtuais e físicos pelo kernel, e pela suposição de que há RAM física no endereço \lstinline{0x80000000}, onde o kernel espera 
ser carregado. Isso funciona com o QEMU, mas em hardware real acaba sendo uma má ideia; hardware real posiciona RAM e dispositivos em endereços físicos imprevisíveis, de modo que (por exemplo) pode 
não haver RAM em \lstinline{0x80000000}, onde o xv6 espera poder armazenar o kernel. Projetos de kernel mais robustos utilizam a tabela de páginas para transformar layouts de memória física 
arbitrários em layouts previsíveis de endereços virtuais do kernel.

O RISC-V oferece proteção no nível dos endereços físicos, mas o xv6 não utiliza esse recurso.

Em máquinas com muita memória pode fazer sentido usar o suporte do RISC-V para ``super páginas''. Páginas pequenas fazem sentido quando a memória física é limitada, pois permitem alocação e troca 
para o disco com granularidade fina. Por exemplo, se um programa usa apenas 8 kilobytes de memória, fornecer-lhe uma super página inteira de 4 megabytes seria um desperdício. Páginas maiores fazem 
sentido em máquinas com muita RAM, e podem reduzir a sobrecarga na manipulação da tabela de páginas.

A ausência, no kernel do xv6, de um alocador semelhante ao {\tt malloc} que possa fornecer memória para objetos pequenos impede o uso de estruturas de dados sofisticadas que exigiriam alocação 
dinâmica. Um kernel mais elaborado provavelmente alocaria blocos pequenos de tamanhos variados, em vez de (como no xv6) apenas blocos de 4096 bytes; um alocador de kernel real precisaria lidar com 
alocações pequenas assim como com grandes.

Alocação de memória é um tema sempre atual, sendo os problemas básicos o uso eficiente da memória limitada e a preparação para requisições futuras desconhecidas~\cite{knuth}. Hoje em dia, as 
pessoas se importam mais com velocidade do que com eficiência de espaço.

%% 
\section{Exercícios}
%% 

\begin{enumerate}
  
\item Analise a árvore de dispositivos do RISC-V para descobrir a quantidade de memória física que o computador possui.

\item Escreva um programa de usuário que cresça seu espaço de endereçamento em um byte chamando \lstinline{sbrk(1)}.
Execute o programa e investigue a tabela de páginas do programa antes da chamada a \lstinline{sbrk} e após a chamada a \lstinline{sbrk}.
Quanto espaço o kernel alocou? O que a PTE da nova memória contém?

\item Modifique o xv6 para utilizar super páginas no kernel.

\item Implementações Unix de \lstinline{exec} tradicionalmente incluem um tratamento especial para scripts de shell. Se o arquivo a ser executado começar com o texto \lstinline{#!}, então a 
primeira linha é interpretada como sendo um programa a ser executado para interpretar o arquivo. Por exemplo, se \lstinline{exec} for chamado para executar \lstinline{myprog} \lstinline{arg1} e a 
primeira linha de \lstinline{myprog} for \lstinline{#!/interp}, então \lstinline{exec} executa \lstinline{/interp} com a linha de comando \lstinline{/interp} \lstinline{myprog} \lstinline{arg1}. 
Implemente suporte a essa convenção no xv6.

\item Implemente randomização do layout do espaço de endereçamento para o kernel.

\end{enumerate}
